\begin{abstract}
The rapid spread of indoor Wi‑Fi has turned everyday spaces into powerful, passive sensing systems—no wearables, no cameras, just the invisible radio waves already around us. From these signals, Channel State Information (CSI) reveals fine-grained distortions caused by human movement, exposing distinctive gait patterns that are extremely hard to fake or imitate. This dissertation harnesses Wi‑Fi CSI for Human Identification (Human Id), rigorously exploring deep learning architectures designed to handle rich, multivariate time‑series CSI data along with auxiliary metadata to push the limits of reliable, device‑free biometric recognition.

We built and tested our models on a rich, publicly available benchmark of CSI traces, captured from 30 subjects performing 5 distinct experiments in both line-of-sight and non-line-of-sight indoor settings. For Gait Pattern, Experiment\#3 is selected having 4 different walking activities. Four cutting-edge deep learning architectures—DenseNet1D, EfficientNet1D-LSTM, MobileNetV3-1D-LSTM, and a custom CSILSTMNet—were all trained within a unified, production-grade MLOps pipeline, leveraging MLflow for end-to-end experiment tracking, systematic hyperparameter optimization, and rock-solid reproducibility.

We systematically pushed the limits of generalization by sweeping learning rates, batch sizes, optimizers, and weight-decay settings. DenseNet1D emerged as the clear winner: with a learning rate of 0.0001, batch size 64, and Adam with 1e-05 weight decay, it delivered the highest validation accuracy (up to 98.33\%) and F1-score (Validation 98.32\%), decisively outperforming all other architectures. MobileNetV3-1D-LSTM and CSILSTMNet also posted strong, competitive results when optimally tuned.

This study exposes how sharply CSI-based identification performance hinges on architectural design and training choices—small tweaks can have big consequences. It further proves that CSI is a powerful, truly passive, and non-intrusive biometric signal, ideal for secure indoor monitoring, with clear potential in high-security facilities, smart buildings, and healthcare settings. Overall, this dissertation delivers a rigorous, end-to-end assessment of deep learning methods for Wi-Fi-based HAR and Human-ID, pinpoints the best-performing architecture, and cements an MLOps-driven framework for fast, scalable, and reproducible experimentation.
\end{abstract}

 
